% ===========================================================================
\section{约束下的熵}\label{sec:entropy-properties}
\warmth{0}\quad\whisper{支撑、条件化与确定性映射都会限制不确定性。}

\begin{strand}
熵不能任意变化。有限字母表固定了它的取值范围，条件化只能消除不确定性，
而确定性映射只能合并结果。取等条件进一步告诉我们：什么时候信息没有损失。
\end{strand}

\block{取值范围与取等条件}

\begin{proposition}[熵的界]\label{prop:entropy-bounds}
若 $X$ 取值于有限字母表 $\X$，则
\[
  0\le \Ent(X)\le\log|\X|.
\]
下界成立当且仅当 $X$ 为常量；上界成立当且仅当 $X$ 在 $\X$ 上均匀分布。
\end{proposition}

\begin{proof}
每一项 $-p_X(x)\log p_X(x)$ 都非负，因此得到下界。等号要求某个结果的
概率为 $1$。

对上界，令 $u$ 为 $\X$ 上的均匀分布。命题~\ref{prop:uniform-kl}
中的均匀参照恒等式给出
\[
  \KL(p_X\Vert u)=\log|\X|-\Ent(X)\ge0.
\]
信息不等式取等当且仅当 $p_X=u$。
\end{proof}

\begin{yourturn}
补全证明上界所用的恒等式：
\[
  \KL(p_X\Vert u)=\log|\X|-\answerin[1.8cm]{\Ent(X)}.
\]
\recall{使用 $\KL(p_X\Vert u)=\log|\X|-\Ent(X)\ge0$。}
\end{yourturn}

\begin{proposition}[条件化降低熵]\label{prop:conditioning-bounds}
对离散随机变量 $X,Y$，
\[
  0\le\Ent(X\mid Y)\le\Ent(X).
\]
此外，
\[
  \Ent(X\mid Y)=\Ent(X)\iff X\text{ 与 }Y\text{ 独立},
\]
而
\[
  \Ent(X\mid Y)=0\iff X=f(Y)\quad\text{几乎处处成立}
\]
其中 $f$ 为某个函数。
\end{proposition}

\begin{proof}
条件熵是非负熵的平均，因此
\[
  \Ent(X\mid Y)=\sum_yp_Y(y)\Ent(X\mid Y=y)\ge0.
\]
它为零当且仅当每个满足 $p_Y(y)>0$ 的条件分布都是点质量；这等价于
$X=f(Y)$ 几乎处处成立。

对上界，命题~\ref{prop:mi-properties} 给出
\[
  \Ent(X)-\Ent(X\mid Y)=\MI(X;Y)\ge0.
\]
等号成立当且仅当 $\MI(X;Y)=0$，也就是 $X,Y$ 独立。
\end{proof}

\begin{yourturn}
配对两个取等条件：
\[
\begin{aligned}
  \Ent(X\mid Y)=\Ent(X)
    &\iff \answerin[3.2cm]{X\perp Y},\\
  \Ent(X\mid Y)=0
    &\iff \answerin[3.2cm]{X=f(Y)\ \text{a.s.}}.
\end{aligned}
\]
\recall{上界取等意味着独立；条件熵为零意味着 $X$ 由 $Y$ 决定。}
\end{yourturn}

\begin{proposition}[熵的链式法则与次可加性]\label{prop:entropy-chain-n}
对离散随机变量 $X_1,\ldots,X_n$，
\[
  \Ent(X_1,\ldots,X_n)
  =\sum_{i=1}^n\Ent(X_i\mid X_1,\ldots,X_{i-1})
  \le\sum_{i=1}^n\Ent(X_i).
\]
不等式取等当且仅当这些随机变量相互独立。
\end{proposition}

\begin{proof}
反复使用二元链式法则
$\Ent(U,V)=\Ent(U)+\Ent(V\mid U)$，再逐项应用
\[
  \Ent(X_i\mid X_1,\ldots,X_{i-1})\le\Ent(X_i).
\]
取等意味着每个 $X_i$ 都与它之前的变量独立，这等价于相互独立。
\end{proof}

\block{确定性映射与碰撞}

\begin{proposition}[函数不会增大熵]\label{prop:function-entropy}
对任意函数 $f$，
\[
  \Ent(f(X))\le\Ent(X).
\]
等号成立当且仅当 $f$ 在 $X$ 的支撑上为单射。
\end{proposition}

\begin{proof}
因为 $f(X)$ 由 $X$ 完全决定，
\[
  \Ent(X,f(X))=\Ent(X)+\Ent(f(X)\mid X)=\Ent(X).
\]
按另一顺序展开：
\[
  \Ent(X,f(X))
  =\Ent(f(X))+\Ent(X\mid f(X))
  \ge\Ent(f(X)).
\]
等号成立当且仅当可以从 $f(X)$ 恢复 $X$，也就是 $f$ 在 $X$ 的支撑上
为单射。
\end{proof}

\begin{yourturn}
哪个非负项度量了 $f$ 丢失的信息？
\[
  \Ent(X)-\Ent(f(X))
  =\answerin[3.4cm]{\Ent(X\mid f(X))}.
\]
\recall{损失的信息是 $\Ent(X\mid f(X))$。}
\end{yourturn}

\begin{proposition}[碰撞概率]\label{prop:collision-probability}
令 $X'$ 为 $X$ 的独立同分布副本，则
\[
  \PP(X=X')\ge 2^{-\Ent(X)}.
\]
等号成立当且仅当 $X$ 在其支撑上均匀分布。
\end{proposition}

\begin{proof}
对凸函数 $t\mapsto2^t$ 使用 Jensen 不等式：
\[
  2^{-\Ent(X)}
  =2^{\E[\log p_X(X)]}
  \le\E\!\left[2^{\log p_X(X)}\right]
  =\E[p_X(X)].
\]
独立同分布给出
\[
  \E[p_X(X)]=\sum_xp_X(x)^2=\PP(X=X').
\]
Jensen 取等当且仅当 $\log p_X(X)$ 在支撑上为常数，也就是 $X$ 在支撑上
均匀分布。
\end{proof}

\begin{yourturn}
补全碰撞恒等式：
\[
  \E[p_X(X)]=\sum_xp_X(x)^2
  =\answerin[2.8cm]{\PP(X=X')}.
\]
\recall{对独立副本 $X'$，$\sum_xp_X(x)^2=\PP(X=X')$。}
\end{yourturn}

\block{混合分布}

\begin{proposition}[熵的凹性]\label{prop:entropy-concavity}
对同一字母表上的概率质量函数 $p,q$ 以及 $0\le\lambda\le1$，
\[
  \Ent\!\bigl(\lambda p+(1-\lambda)q\bigr)
  \ge\lambda\Ent(p)+(1-\lambda)\Ent(q).
\]
\end{proposition}

\begin{proof}
函数 $g(t)=-t\log t$ 在 $[0,1]$ 上为凹函数。因此对每个 $x$，
\[
  g\!\bigl(\lambda p(x)+(1-\lambda)q(x)\bigr)
  \ge\lambda g(p(x))+(1-\lambda)g(q(x)).
\]
对 $x$ 求和即得结论。
\end{proof}

\begin{yourturn}
写出熵的凹性背后的一元函数：
\[
  g(t)=\answerin[2.4cm]{-t\log t}.
\]
\recall{熵是凹函数项 $g(t)=-t\log t$ 的和。}
\end{yourturn}
